\chapter{Bullet 3: zero-shot Flan-T5}
\label{ch:flan}

\begin{cvline}
Benchmarked zero-shot Flan-T5 models, achieving 0.87 ROUGE-L on Flan-T5-large and demonstrating a 6.6\% lift via prompt syntax tuning
\end{cvline}

%=====================================================================
\section{Words you need}

\begin{definition}[Instruction tuning]
\dfn{Instruction tuning}: After pretraining, train the model on very many tasks, each written as an instruction plus an answer (``Translate this to German: \dots''). The model learns to follow instructions it has never seen. Flan-T5 is T5 after instruction tuning on over a thousand tasks.
\end{definition}

\begin{definition}[Zero-shot]
\dfn{Zero-shot}: Using a model on a task with no examples of that task and no extra training. Here Flan-T5 gets only an instruction and the article. Not one of the 13.4k pairs is used.
\end{definition}

\begin{definition}[Prompt syntax tuning]
\dfn{Prompt syntax tuning}: Trying different wordings of the instruction and keeping the one that scores best. The model's weights never change; only the text in front of the article does. You tried two:
\begin{enumerate}
  \item ``Summarize the following text into a title: \{text\}''
  \item ``Generate a headline for this article: \{text\}''
\end{enumerate}
\end{definition}

\begin{definition}[Model size]
\dfn{Model size} is the number of weights. Flan-T5-base has about 250 million; Flan-T5-large about 780 million, roughly 3 times as many and 13 times T5-small's 60 million.
\end{definition}

\begin{definition}[Relative lift]
The \dfn{relative lift} from score $a$ to score $b$ is $(b - a)/a$: the gain as a fraction of where you started. It differs from the plain difference $b - a$.
\end{definition}

\begin{example}[Your 6.6\%]
Flan-T5-large's ROUGE-L is 0.8207 with prompt 1 and 0.8745 with prompt 2.
\begin{steps}
  \item Plain difference: $0.8745 - 0.8207 = 0.0538$.
  \item Relative lift: $0.0538 / 0.8207 = 0.0656 \approx 6.6\%$.
\end{steps}
\end{example}

\begin{definition}[Resolution of a test set]
The \dfn{resolution} is the smallest change a test set can show. ROUGE-L here is an average over 100 articles, each scoring between 0 and 1, so one article changing its score completely moves the average by $1/100 = 0.01$. Gaps much under 0.01 can't be told apart from noise.
\end{definition}

\begin{sayit}[Say it in 20 seconds]
I also tried Flan-T5, an instruction-tuned T5, with no training at all, just an instruction in front of the article. Two sizes, base and large, and two prompt wordings. Flan-T5-large reached 0.87 ROUGE-L, level with fine-tuned T5-small. Rewording the prompt from ``summarize into a title'' to ``generate a headline'' lifted it from 0.82 to 0.87, a 6.6\% relative lift with zero training.
\end{sayit}

%=====================================================================
\section{The picture}

\begin{sayit}[Board Q · Fine-tuning against prompting]
\begin{pic}
\begin{tikzpicture}[>={Latex[length=2mm]},font=\sffamily\small,
  box/.style={draw=#1L,fill=#1,rounded corners=2pt,align=center,minimum height=9mm,inner sep=3pt}]
\node[font=\sffamily\small\bfseries] at (2.2,1.3) {T5-small: fine-tuned};
\node[box=bdq] (a) at (0,0) {13.4k pairs};
\node[box=bdv] (b) at (3.6,0) {weights change};
\draw[->] (a) -- (b);
\node[font=\sffamily\small\bfseries] at (8.6,1.3) {Flan-T5: zero-shot};
\node[box=bdy] (c) at (7,0) {an instruction};
\node[box=bdg] (d) at (10.4,0) {weights fixed};
\draw[->] (c) -- (d);
\node[font=\sffamily\scriptsize] at (2.2,-1) {0.86 ROUGE-L}; \node[font=\sffamily\scriptsize] at (8.6,-1) {0.87 ROUGE-L (large, prompt 2)};
\end{tikzpicture}
\end{pic}
Say: ``One model learned the task from 13.4k examples; the other was only told what to do. They ended level.''
\end{sayit}

\begin{sayit}[Board R · The 2 $\times$ 2 grid (ROUGE-L)]
\begin{pic}
\begin{tikzpicture}[font=\sffamily\small]
\node at (3,1.3) {prompt 1: ``summarize \dots into a title''}; \node at (8.6,1.3) {prompt 2: ``generate a headline''};
\node[anchor=east] at (0.4,0.4) {Flan-T5-base}; \node[anchor=east] at (0.4,-0.8) {Flan-T5-large};
\draw[fill=defc!20] (0.6,0) rectangle (5.4,0.8); \node at (3,0.4) {0.6808};
\draw[fill=defc!12] (6.2,0) rectangle (11,0.8); \node at (8.6,0.4) {0.6317};
\draw[fill=thmc!20] (0.6,-1.2) rectangle (5.4,-0.4); \node at (3,-0.8) {0.8207};
\draw[fill=thmc!40] (6.2,-1.2) rectangle (11,-0.4); \node at (8.6,-0.8) {\textbf{0.8745}};
\end{tikzpicture}
\end{pic}
Say: ``Size matters most: large beats base by 0.14 to 0.24. The prompt matters too, and the best wording isn't the same for both sizes.''
\end{sayit}

\begin{sayit}[Board S · Where the 6.6\% comes from]
\begin{pic}
\begin{tikzpicture}[font=\sffamily\small,x=11cm]
\fill[thmc!20,draw=thmc] (0,0) rectangle (0.8207,0.5); \node[anchor=west] at (0.01,0.25) {large, prompt 1: 0.8207};
\fill[thmc!45,draw=thmc] (0,-0.75) rectangle (0.8745,-0.25); \node[anchor=west] at (0.01,-0.5) {large, prompt 2: 0.8745};
\draw[<->,trapc,thick] (0.8207,-0.9) -- (0.8745,-0.9);
\node[text=trapc,anchor=north,font=\sffamily\scriptsize] at (0.8476,-0.95) {0.0538: $0.0538/0.8207 = 6.6\%$};
\end{tikzpicture}
\end{pic}
Say: ``Same model, same weights; only the instruction changed. 0.054 on 0.821 is a 6.6\% relative lift: about five articles' worth.''
\end{sayit}

\begin{sayit}[Board T · Every model on one line (ROUGE-L)]
\begin{pic}
\begin{tikzpicture}[font=\sffamily\scriptsize,x=12cm]
\draw[->] (0.2,0) -- (0.95,0);
\foreach \x in {0.2,0.3,0.4,0.5,0.6,0.7,0.8,0.9} {\draw (\x,0.06) -- (\x,-0.06) node[below]{\x};}
\foreach \v/\n/\c/\h/\a in {0.2469/basic GRU/prec/0.5/south,0.4985/GRU + GloVe/defc/0.5/south,0.6808/Flan-base/exc/0.5/south,0.8207/Flan-large p1/thmc/0.8/south east,0.8626/T5-small/trapc/1.3/south west,0.8745/Flan-large p2/thmc/1.9/south west} {
  \draw[\c,thick] (\v,0) -- (\v,\h); \node[\c,anchor=\a] at (\v,\h) {\n};}
\end{tikzpicture}
\end{pic}
Say: ``The three families sit in three bands: RNNs at 0.25--0.50, Transformers at 0.68--0.87. T5-small and Flan-T5-large are about one article apart.''
\end{sayit}

%=====================================================================
\section{How it works}
\begin{enumerate}
  \item Load \code{google/flan-t5-base} and \code{google/flan-t5-large} with their pretrained weights; no training.
  \item For each test article, put one of the two instructions in front of the raw text.
  \item Generate with beam search (width 5, up to 50 tokens) and score with ROUGE. (These runs were on CPU, which is why they were slow.)
  \item Compare the four combinations: 2 sizes $\times$ 2 prompts.
\end{enumerate}

%=====================================================================
\section{The numbers: 0.87 and 6.6\%}

\begin{measured}[The number]
\begin{itemize}
  \item \textbf{0.87 ROUGE-L:} Flan-T5-large, prompt 2, scored 0.8745.
  \item \textbf{6.6\% lift:} $(0.8745 - 0.8207)/0.8207$; about 5 articles' worth, so a real effect.
  \item \textbf{Level with fine-tuning:} $0.8745 - 0.8626 = 0.012$ above T5-small, about one article: call it a tie. (On ROUGE-2 T5-small is slightly ahead.)
  \item \textbf{The catch:} the Flan scores used ROUGE without stemming, T5's with; each is one run on 100 articles.
\end{itemize}
\end{measured}

%=====================================================================
\section{What they will ask}

\begin{iq}[\core]{How can a model with no training match a fine-tuned one?}
Scale and instruction tuning. Flan-T5-large is 13 times the size of T5-small and has been trained to follow instructions across over a thousand tasks, including summarization and headlines. So ``write a headline for this'' is a task it effectively knows. Fine-tuning gives T5-small the task; Flan-T5-large brings it from scale.
\end{iq}

\begin{iq}[\core]{Why would the prompt wording matter?}
The model learned from instructions written in particular ways. ``Headline'' matches a familiar task, short and title-like; ``summarize into a title'' can pull it towards a longer, summary-like answer. The weights are fixed, so the wording is the only lever, and it moved ROUGE-L by 0.054.
\end{iq}

\begin{iq}[\core]{Is 0.87 really better than 0.86?}
No: 0.012 is about one article's worth, within noise for a single run on 100 articles, and ROUGE-2 points the other way. The fair claim is that zero-shot Flan-T5-large matched fine-tuned T5-small.
\end{iq}

\begin{iq}{Isn't trying prompts on the test set a form of tuning?}
Yes, a mild one: choosing the better of two prompts by test score flatters the winner slightly. The clean way is to choose on the validation articles and report the test score once.
\end{iq}

\begin{iq}{Why did base and large prefer different prompts?}
Smaller models follow instructions less reliably, so which wording works is less predictable. It's one run each, so I wouldn't read much into the flip; the size effect is the robust finding.
\end{iq}

\begin{iq}{What would you do next with Flan-T5?}
Fine-tune it: Flan-T5-base or -large on the 13.4k pairs, with LoRA to keep it cheap, which trains only small extra matrices instead of every weight. That combines scale with task data.
\end{iq}

\begin{bugfix}[Don't say]
\begin{itemize}
  \item ``A 6.6-point lift.'' It's 6.6\% relative; 0.054 absolute.
  \item ``Flan-T5 beat fine-tuning.'' It matched it, within one article.
  \item ``Flan-T5 was fine-tuned.'' It was zero-shot: no training.
\end{itemize}
\end{bugfix}
